\section{Study 1: The Punitiveness Ecosystem}

Study~1 provides a comprehensive correlational map of the psychological factors associated with punitiveness, designed to enable side-by-side comparison of the strengths of the normatively appropriate and inappropriate clusters. Specifically, we test the relative strength of the association between punitiveness and the crime concerns cluster against the association between punitiveness and the hostile aggression cluster.

\subsection{Method}

This study was preregistered on the Open Science Framework prior to data collection (\url{https://osf.io/kr7y2/}). All data, materials, analysis code, and an interactive results website are publicly available at \url{https://dgk-law-and-cognition-lab.github.io/Punishment-Punitiveness/}. We report how we determined our sample size, all data exclusions, all manipulations, and all measures \citep{simmons2012false}. The quantitative analyses were conducted in R \citep{Rcore2024}.

\subsubsection{Participants}

A sample of 496 U.S. adults was recruited through Prolific, using Prolific's politically representative US sampling feature, which applies quota balancing on age, sex, ethnicity, and political affiliation. The preregistered target was 480, based on the sample size of a preliminary version of the study \citep{SimonMelnikoff2025}. From an initial pool of 538 respondents, we excluded those who did not complete the survey ($n = 8$) and those who failed either of two embedded attention checks ($n = 34$). Participants were compensated at Prolific's standard rate. The study was approved by the Institutional Review Board at the University of Southern California.

The final sample had a mean age of 46.2 years ($SD = 16.2$, range: 19--85). Gender composition was 49.2\% male, 50.2\% female, and 0.6\% other. Racial and ethnic composition was 65.3\% White, 13.3\% Black, 10.7\% Hispanic/Latino, 7.9\% Asian, and 2.8\% other. Political identification was evenly distributed, with 38.3\% liberal, 23.4\% moderate, and 38.3\% conservative.

\subsubsection{Design and Procedure}

The study employed a correlational design. All participants completed the same battery of measures in a fixed order, with one exception: for the sentencing task, each participant was randomly assigned to one of three criminal vignettes, designed to enhance the generalizability of the findings across crime types. The survey began with punitiveness attitude measures, followed by crime concern measures, hostile aggression and emotion measures, personality and ideology measures, the randomly assigned criminal vignette with a sentencing decision, an open-ended sentencing justification (the source of data for Study~2), and demographics (for the wording of all measures, see Supplementary Section~\ref{sec:supp_materials}). In all, the survey included sixty-nine measures, presented in twelve blocks. The blocks were presented in a fixed order, but the measures within the blocks were randomly presented. All quantitative measures used 7-point Likert scales (1 = \textit{Strongly disagree} to 7 = \textit{Strongly agree}) unless otherwise noted. Reliability estimates for all scales are reported in Table~\ref{tab:reliability}.

The sentencing measure, which constituted part of the punitiveness construct, was presented towards the end of the survey so as not to influence the other measures. To increase the generalizability of the sentencing measure, participants were presented with one of three randomized vignettes in which a man was convicted for second-degree murder (Stranger Felony-Murder, $n = 168$; Domestic Violence, $n = 176$; Organized Crime, $n = 152$; full vignette text in Supplementary Section~\ref{sec:supp_vignettes}). Participants were informed that judges typically gave sentences between twenty and thirty years for that crime, and were then asked to recommend a sentence on a scale of 0--50 years. Following the sentencing recommendation, participants were presented with an open-ended prompt: ``In your own words, please explain why you recommended this sentence for Darryl. What was your reasoning?'' A follow-up prompt encouraged elaboration. These open-ended responses are analyzed in Study~2.

\subsubsection{Measures}

\paragraph{Punitiveness.} The primary dependent variable was a standardized composite capturing three distinct operationalizations of punitive attitudes. \textit{Penological attitudes} were assessed with four items measuring endorsement of more punishment (2 items, e.g., ``The solution to crime is more punishment'') and rejection of the parsimony principle (2 items, e.g., ``The state should not punish more than is absolutely necessary''). \textit{Policy preferences} were assessed with four items tapping support for three-strikes laws (2 items, e.g., ``People convicted three times for a felony should be imprisoned for life''), life without parole (1 item: ``Convicted murderers should never be allowed out of prison''), and the death penalty (1 item: ``How much do you support or oppose the use of the death penalty to punish people convicted of first-degree murder?''). These eight items formed a reliable composite ($\alpha = .84$). The \textit{sentencing decision}---the number of years assigned to the defendant---was $z$-scored within each vignette condition to account for cross-vignette differences in mean sentence length, and then averaged with the $z$-scored 8-item attitude composite to produce the overall punitiveness aggregate. When the standardized sentencing decision was included alongside the eight attitude items, internal consistency remained excellent ($\alpha = .84$).

\paragraph{Correlate Measures.} The correlate measures were organized into four thematic clusters. The \textit{crime concerns} cluster comprised perceived crime rates (2 items, e.g., ``Do you feel that violent crime just keeps rising?''; $\alpha = .85$) and fear of crime (3 items, e.g., ``How safe do you feel walking at night in the vicinity of your home?''; $\alpha = .78$). The \textit{emotional reactions} cluster comprised hatred toward criminals (3 items, e.g., ``I feel hatred towards criminals''; $\alpha = .78$) and anger toward criminals (2 items, e.g., ``I feel anger towards criminals''; $\alpha = .83$).

The \textit{hostile aggression} cluster comprised six constructs: support for social exclusion of criminals (3 items, e.g., ``Having chosen to break the law, criminals do not deserve to enjoy the benefits that society provides its members''; $\alpha = .76$); endorsement of degradation (3 items, e.g., ``Punishment should humiliate the criminal''; $\alpha = .67$); endorsement of inflicting suffering (2 items, e.g., ``Infliction of suffering is a permissible goal of criminal punishment''; $\alpha = .79$); tolerance of prison violence (2 items, e.g., ``Sexual assault of prisoners is a serious societal problem,'' reverse-coded; $\alpha = .56$); endorsement of harsh prison conditions (3 items, e.g., ``The conditions in prisons should be extremely harsh''; $\alpha = .82$); and endorsement of revenge (3 items, e.g., ``Society has the right to take revenge on criminal offenders''; $\alpha = .67$). The composite of all six constructs showed excellent reliability ($\alpha = .92$).

The \textit{personality and ideology} cluster comprised right-wing authoritarianism (RWA; 5 items, e.g., ``What our country needs most is discipline, with everyone following our leaders in unity''; $\alpha = .87$), social dominance orientation (SDO; 8 items, e.g., ``An ideal society requires some groups to be on top and others to be on the bottom''; $\alpha = .92$), trait vengefulness (5 items, e.g., ``There is nothing wrong in getting back at someone who has hurt you''; $\alpha = .88$), violence proneness (4 items, e.g., ``Physical punishment is necessary for raising children''; $\alpha = .73$), racial resentment (4 items, e.g., ``Irish, Italian, Jewish, and many other minorities overcame prejudice and worked their way up. Blacks should do the same without any special favors''; $\alpha = .90$), and blood sports viewership (4 items, e.g., ``How much do you like watching Mixed Martial Arts (UFC) on TV?''; $\alpha = .84$).

Two additional constructs, organized into a \textit{process violations} cluster, were collected for exploratory purposes: tolerance of due process violations (3 items, e.g., ``The legal system's loopholes and technicalities make it too difficult to convict criminals for their crimes''; $\alpha = .60$) and willingness to convict on uncertain evidence (4 items, e.g., ``If police officers know that a defendant is guilty, it is okay for them to lie when testifying in court to help obtain a conviction''; $\alpha = .67$). This cluster was dropped from the confirmatory hypotheses but is reported in the exploratory analyses.

Personality measures were borrowed verbatim from established scales: right-wing authoritarianism \citep{BizumicDuckitt2018}, social dominance orientation \citep{HoEtAl2015}, and racial resentment \citep{KinderSanders1996}. Other measures were borrowed or adapted from \citet{Cullen1985}, \citet{ChiricosWelchGertz2004}, \citet{LecciMyers2008}, \citet{KingMaruna2009}, \citet{TylerBoeckmann1997}, and \citet{Wozniak2016}. To the best of our knowledge, the following constructs have not been examined in relation to punitiveness in prior work: willingness to convict on uncertain evidence, hatred, infliction of suffering, tolerance of prison violence, violence proneness, and blood sports viewership. The items used to measure these constructs were formulated for the present study.

\subsubsection{Analytic Strategy}

The correlational analyses used Pearson correlations between punitiveness and all correlate constructs. To test whether hostile aggression outpredicts crime concerns (H2), we applied Steiger's $Z$-test for comparing dependent correlations \citep{steiger1980tests} and computed bootstrapped 95\% confidence intervals for the difference between dependent correlations (10{,}000 iterations, BCa intervals). False discovery rate (FDR) correction was applied for the 18 simultaneous construct-level tests of H1.

\subsection{Results}

\subsubsection{Preliminary Analyses}

The final sample comprised 496 U.S. adults randomly assigned to one of three vignettes (Stranger Felony-Murder $n = 168$, Domestic Violence $n = 176$, Organized Crime $n = 152$). The punitiveness composite ($M = 0.00$, $SD = 0.82$) was centered at zero by construction. All primary cluster composites demonstrated good to excellent internal consistency (range: $\alpha = .78$ to $.92$). Two individual constructs---parsimony ($\alpha = .48$) and prison violence tolerance ($\alpha = .56$)---fell below conventional thresholds; results involving these measures are interpreted with caution.

\paragraph{Sentencing Distribution.} Recommended sentences departed substantially from the 25-year reference point conveyed in the instructions. Pooled across vignettes, the mean sentence was 35.2 years ($SD = 12.5$, $\mathit{Mdn} = 35$), 10.2 years above the anchor, $t(495) = 18.14$, $p < .001$, $d = 0.81$. Per-vignette means were 34.4 years (Vignette~A: $d = 0.75$), 37.5 years (Vignette~B: $d = 0.99$), and 33.4 years (Vignette~C: $d = 0.70$); all three differed reliably from 25 years (all $p < .001$). The deviation held across the political spectrum: liberals ($M = 33.0$, $d = 0.62$), moderates ($M = 37.1$, $d = 0.94$), and conservatives ($M = 36.2$, $d = 0.97$) each sentenced significantly above the anchor (all $p < .001$). A one-way ANOVA confirmed a small but reliable vignette effect, $F(2, 493) = 4.84$, $p = .008$, $\eta^2 = .019$, with Vignette~B yielding harsher sentences than Vignette~C (Tukey $p = .009$). The right-skewed distributions, with 32.3\% of participants assigning the maximum 50 years, are displayed in Figure~\ref{fig:sentence_anchor}. Notably, participants gravitated toward round-number values: 86\% of all responses concentrated on just six values (50, 30, 25, 40, 20, and 35 years), reflecting a tendency toward round-number anchoring rather than fine-grained calibration along the 0--50 scale. Because the sentencing decision was $z$-scored within vignette before entering the punitiveness aggregate, this between-vignette variation does not bias the correlational analyses that follow.

\subsubsection{Hypothesis 1: Punitiveness Correlates With All Measures}

Hypothesis~1 predicted that punitiveness would be positively correlated with all correlate measures. This prediction was strongly supported. All 18 construct-level correlations were positive, and 16 of 18 were significant after FDR correction (Table~\ref{tab:H1_correlations}). The strongest correlates were endorsement of harsh prison conditions ($r = .56$), social exclusion ($r = .54$), hatred of criminals ($r = .53$), infliction of suffering ($r = .52$), right-wing authoritarianism ($r = .50$), and racial resentment ($r = .46$). The two exceptions were trait vengefulness ($r = .12$, $p_\textit{FDR} = .085$) and blood sports viewership ($r = .11$, $p_\textit{FDR} = .080$), which were positive but did not survive correction. By contrast, the crime concern measures---perceived crime rates ($r = .40$) and fear of crime ($r = .18$)---correlated substantially less strongly with punitiveness than did the hostile aggression and ideology constructs.

The macro-level pattern across clusters (Figure~\ref{fig:heatmap_cluster}) reinforces the construct-level findings. The four thematic clusters all correlated positively with punitiveness, but the magnitude varied by an order of psychological proximity to harsh punishment: hostile aggression ($r = .59$, 95\% CI $[.53, .64]$) $>$ emotions ($r = .53$, $[.47, .60]$) $>$ personality and ideology ($r = .47$, $[.40, .54]$) $>$ crime concerns ($r = .32$, $[.24, .40]$).

\subsubsection{Hypothesis 2: Hostile Aggression Outpredicts Crime Concerns}

The central confirmatory test compared the strengths of the punitiveness--hostile aggression and punitiveness--crime concerns associations. The hostile aggression cluster correlated with punitiveness at $r = .59$, whereas the crime concerns cluster correlated at $r = .32$. Steiger's $Z$-test for dependent correlations confirmed that this difference was reliable: $Z = 6.26$, $p < .001$. A bootstrapped 95\% confidence interval for the difference between the two dependent correlations excluded zero, $\Delta r = .264$, BCa CI $[.174, .367]$ (Figure~\ref{fig:bootstrap}).

The same pattern held when each individual punitiveness measure was treated as the dependent variable. Hostile aggression significantly outpredicted crime concerns for five of six measures: Punish More ($\Delta r = .34$, $Z = 8.23$, $p < .001$), Parsimony ($\Delta r = .31$, $Z = 6.70$, $p < .001$), Death Penalty ($\Delta r = .32$, $Z = 7.23$, $p < .001$), Three Strikes ($\Delta r = .22$, $Z = 5.06$, $p < .001$), and Life Without Parole ($\Delta r = .15$, $Z = 3.27$, $p = .001$). The one exception was the sentencing decision itself ($\Delta r = .08$, $Z = 1.60$, $p = .11$).

\paragraph{Tautology Sensitivity Analysis.} A potential concern is that some hostile aggression constructs overlap conceptually with punitiveness, which could artifactually inflate the comparison. We addressed this concern through a graduated sensitivity analysis at three stringency levels (Table~\ref{tab:steiger}; Figure~\ref{fig:bootstrap}). At the \textit{original} level (all six hostile aggression constructs: exclusion, degradation, suffering, prison violence, harsh conditions, revenge), the advantage was $\Delta r = .264$, $Z = 6.26$, $p < .001$, BCa CI $[.174, .367]$. At a \textit{conservative} level that retained the four constructs least prone to tautology (dropping suffering and harsh conditions), the advantage held: $\Delta r = .226$, $Z = 5.07$, $p < .001$, BCa CI $[.134, .331]$. At the \textit{strictest} level---retaining only social exclusion and revenge, the two constructs with the most distinct content from the punitiveness measures---the advantage remained substantial: $\Delta r = .221$, $Z = 5.05$, $p < .001$, BCa CI $[.125, .323]$. All three bootstrap intervals excluded zero. As an additional robustness check, dropping the low-reliability parsimony items from the punitiveness composite produced virtually identical results ($r = .58$ vs.\ $r = .33$, $Z = 5.87$, $p < .001$). The hostile aggression advantage is therefore not an artifact of construct overlap with the dependent variable.

\subsubsection{Hypothesis 3: Cohesive Punitiveness Mindset}

Hypothesis~3 predicted that the correlate constructs would themselves be positively intercorrelated, indicating a cohesive psychological mindset rather than a heterogeneous set of independent factors. The 18 correlate constructs yielded 153 pairwise correlations. All 153 (100\%) were positive, ranging from $r = .02$ to $r = .72$, and 144 (94.1\%) were significant at $p < .05$. The hostile aggression constructs formed an especially tight subcluster (within-cluster $r$s ranging from $.38$ to $.71$). The punitiveness--correlate correlations were also stable across the three vignette conditions (median range across vignettes $= .14$), with no construct reversing direction in any vignette.

\subsubsection{Demographic Moderation}

To test whether the punitiveness ecosystem operates differently across demographic groups, we ran a comprehensive set of moderation analyses. For each of four key correlates (Hostile Aggression, Crime Concerns, RWA, Racial Resentment), we tested interactions with six demographic moderators (age, gender, education, race [white vs.\ non-white], race [7-level], and political orientation), yielding 24 interaction tests for the punitiveness aggregate (Table~\ref{tab:demographic_moderation_punit}) and 24 for the (within-vignette $z$-scored) sentencing decision (Table~\ref{tab:demographic_moderation_sent}).

After FDR correction for the 48 simultaneous tests, only one interaction reached significance: Crime Concerns $\times$ Age predicting the punitiveness aggregate ($b = 0.16$, $p_\textit{FDR} = .000061$), reflecting a slightly stronger crime-concern--punitiveness link among older participants. The remaining 47 interactions (including all four hostile aggression interactions, all four RWA interactions, all four racial resentment interactions, and the political orientation $\times$ hostile aggression interaction, $b = -0.022$, $p_\textit{raw} = .074$, $p_\textit{FDR} = .444$) were non-significant after correction. The punitiveness ecosystem documented above is therefore largely demographic-invariant: hostile aggression, ideology, and racial resentment correlate with punitiveness at comparable strengths across age groups, genders, racial categories, education levels, and the liberal--conservative spectrum.

\subsubsection{Confirmatory Factor Analysis}

To examine the latent factor structure underlying the four confirmatory clusters, we fit a CFA model in which the 16 confirmatory correlate composites loaded on four latent factors corresponding to Crime Concerns, Emotional Reactions, Hostile Aggression, and Personality \& Ideology. All 16 standardized loadings were positive and significant ($\lambda$s ranging from $.37$ for blood sports to $.91$ for hatred; all $p < .001$). Inter-factor correlations were uniformly large, with the strongest links between Emotional Reactions and Hostile Aggression ($r = .80$) and between Hostile Aggression and Personality \& Ideology ($r = .78$). Global model fit, however, fell short of conventional thresholds, $\chi^2(98) = 699.10$, $p < .001$, CFI $= .858$, TLI $= .826$, RMSEA $= .111$ [.104, .119], SRMR $= .069$. The imperfect fit indicates that the four-cluster organization is best understood as a conceptual scaffold for interpreting the correlations among constructs rather than a strict latent measurement model. The analyses reported above operate at the level of the observed cluster composites, which makes the substantive findings robust to this measurement-model limitation.

\subsection{Discussion}

Study~1 reveals that punitiveness is embedded in a dense, cohesive psychological ecosystem. The most prosocially justifiable correlate of punitiveness---concern about crime---is also its weakest predictor ($r = .32$), especially when compared with the hostile aggression cluster ($r = .59$). Notably high correlations with punitiveness were also observed for hatred of criminals ($r = .53$), infliction of suffering ($r = .52$), endorsement of revenge ($r = .43$), authoritarianism ($r = .50$), and racial resentment ($r = .46$). This pattern survived stringent tautology sensitivity tests, held across all three crime vignettes, and showed essentially no demographic moderation.

These findings converge with and extend prior research linking punitiveness to a variety of psychological phenomena which, as discussed above, has been gleaned primarily through a large number of dispersed studies \citep{GerberJackson2013, GoldbergLernerTetlock1999, Ho2002, Jackson2019}. What the present study adds is the capability to observe the relative magnitude of these associations. When hostile aggression, crime concerns, emotional reactions, and personality factors are examined side by side, we see clear evidence that it is the aggressive, hateful, vengeful, and hierarchical dispositions that show the strongest ties to punitive attitudes---not the concern about crime that would be expected under the prosocial account.

A critical question remains, however: if punitiveness is driven primarily by hostile motivations rather than prosocial concerns, why do people so reliably invoke prosocial justifications when explaining their punishment preferences? Study~2 addresses this question directly by analyzing the language people use to justify their sentencing decisions.